% Source: OpenAI2026AreaLaw, Lemma 7.2, 06-transport.tex, lines 144--172.
% Manuscript revision: adc7f1241b42e322a6451854ab7e4b4c146bf78a.
\section{Intertwining the recursive maps of a mean tree}
\label{sec:mean_tree_derivative_intertwine}

Let $T$ be a weighted binary tree as in Definition~\ref{def:mean_tree},
with labels in a set $I$ and positive definite inputs $(A_i)_{i\in I}$.
Write $M_T(A)$ for its root. For a label $j$, define the complex-linear
map $D_{T,j}^A$ recursively. At a leaf labelled $i$, it is the identity
map if $i=j$ and the zero map otherwise. At a vertex with parameter
$p\in[0,1]$ and subtrees $l,r$, it is
\begin{align}
  D_{T,j}^A[H]
    &=D_1F_p\bigl(M_l(A),M_r(A)\bigr)[D_{l,j}^A[H]]
      +D_2F_p\bigl(M_l(A),M_r(A)\bigr)[D_{r,j}^A[H]].
\end{align}
Here $D_1F_p,D_2F_p$ are the complex-linear maps of
Theorem~\ref{thm:geometric_mean_derivative_intertwine}.
Repeated labels contribute through both subtrees whenever they occur in both.

Let $w_{T,j}$ be the terminal weight of label $j$, summing over all its
occurrences. The associated normalized map is
\begin{align}
  \Phi_{T,j}^A[H]
    &=w_{T,j}^{-1}M_T(A)^{-1/2}
      D_{T,j}^A[A_j^{1/2}HA_j^{1/2}]M_T(A)^{-1/2},
\end{align}
with the convention $0^{-1}=0$. Thus this definition also applies when
$w_{T,j}=0$. For distinct labels it is the normalized map appearing in
Theorem~\ref{thm:mean_tree_leaf_maps}. The identities below concern these
recursive maps, including repeated labels, and are auxiliary to
Lemma~7.2 of~\cite{OpenAI2026AreaLaw}.

\begin{theorem}[Intertwining the recursive derivative and normalized map]
  \label{thm:mean_tree_derivative_intertwine}
  \lean{Matrix.MeanTree.derivLabel_intertwine,
    Matrix.MeanTree.leafMap_intertwine}
  \leanok
  Let $(A_i)_{i\in I}$ and $(B_i)_{i\in I}$ be families of positive definite
  matrices on two finite-dimensional complex Hilbert spaces.
  Suppose that a linear map $J$ from the second space to the first satisfies
  $A_iJ=JB_i$ for every $i\in I$. For the same tree $T$, every label $j$
  and every matrix $H$ on the second space satisfy
  \begin{align}
    D_{T,j}^A[JHJ^*]&=J D_{T,j}^B[H]J^*, \\
    \Phi_{T,j}^A[JHJ^*]&=J\Phi_{T,j}^B[H]J^*.
  \end{align}
  No distinctness of labels, positivity of $w_{T,j}$, Hermitian assumption
  on $H$, or isometry assumption on $J$ is required.
\end{theorem}

\begin{proof}
  \leanok
  \uses{def:mean_tree, thm:mean_tree_properties, thm:mean_tree_intertwine,
    thm:geometric_mean_derivative_intertwine, thm:real_power_intertwine,
    thm:matrix_sandwich_intertwine}
  Induct on $T$ for the identity involving $D_{T,j}$.
  At a leaf, the identity and zero maps both satisfy the assertion.
  At a vertex, Theorem~\ref{thm:mean_tree_intertwine} intertwines the
  positive definite subtree roots. Apply
  Theorem~\ref{thm:geometric_mean_derivative_intertwine} to each term
  in the recursion, using the two induction hypotheses, and add the
  resulting identities.

  For the normalized map, Theorem~\ref{thm:real_power_intertwine} gives
  \begin{align}
    A_j^{1/2}J&=JB_j^{1/2}, \\
    M_T(A)^{-1/2}J&=JM_T(B)^{-1/2}.
  \end{align}
  These powers are Hermitian. Apply
  Theorem~\ref{thm:matrix_sandwich_intertwine} to the inner square-root
  sandwich and the outer inverse-square-root sandwich, with the
  derivative intertwining identity between them. Multiplication by the
  common scalar $w_{T,j}^{-1}$ proves the normalized identity, including
  $w_{T,j}=0$.
\end{proof}

\begin{theorem}[Compression of the recursive derivative and normalized map]
  \label{thm:mean_tree_derivative_compress}
  \lean{Matrix.MeanTree.derivLabel_compress,
    Matrix.MeanTree.leafMap_compress}
  \leanok
  Under the hypotheses of Theorem~\ref{thm:mean_tree_derivative_intertwine},
  every matrix $X$ on the first space satisfies
  \begin{align}
    J^*D_{T,j}^A[X]J&=D_{T,j}^B[J^*XJ], \\
    J^*\Phi_{T,j}^A[X]J&=\Phi_{T,j}^B[J^*XJ].
  \end{align}
\end{theorem}

\begin{proof}
  \leanok
  \uses{thm:matrix_adjoint_intertwine, thm:mean_tree_derivative_intertwine}
  The inputs are Hermitian, so
  Theorem~\ref{thm:matrix_adjoint_intertwine} gives $B_iJ^*=J^*A_i$
  for every label $i$. Apply
  Theorem~\ref{thm:mean_tree_derivative_intertwine} with the two input
  families exchanged and intertwiner $J^*$, using $(J^*)^*=J$.
\end{proof}
